%!TEX root = ../hutbthesis_main.tex
\chapter{主要符号、缩写说明、方法对比表格}

\begin{longtable}[]{@{}
  >{\raggedright\arraybackslash}p{(\columnwidth - 2\tabcolsep) * \real{0.1551}}
  >{\raggedright\arraybackslash}p{(\columnwidth - 2\tabcolsep) * \real{0.8449}}@{}}
\toprule\noalign{}
\begin{minipage}[b]{\linewidth}\raggedright
\textbf{缩写/符号}
\end{minipage} & \begin{minipage}[b]{\linewidth}\raggedright
\textbf{含义}
\end{minipage} \\
\midrule\noalign{}
\endhead
\bottomrule\noalign{}
\endlastfoot
ROS & Robot Operating System，机器人操作系统 \\
Gazebo & 机器人三维物理仿真平台 \\
URDF & Unified Robot Description Format，机器人统一描述格式 \\
A-SMGCS & 高级场面活动引导与控制系统 \\
LiDAR & 激光雷达 \\
\(v\) & 车辆纵向速度 \\
\(\omega\) & 车辆偏航角速度 \\
\(\delta_{f}\) & 前桥转角 \\
\(\delta_{r}\) & 后桥转角 \\
\(e_{x}\) & 纵向误差 \\
\(e_{y}\) & 横向误差 \\
\(e_{\theta}\) & 姿态误差 \\
\end{longtable}

方法对比表格

\begin{longtable}[]{@{}
  >{\raggedright\arraybackslash}p{(\columnwidth - 8\tabcolsep) * \real{0.1297}}
  >{\raggedright\arraybackslash}p{(\columnwidth - 8\tabcolsep) * \real{0.2442}}
  >{\raggedright\arraybackslash}p{(\columnwidth - 8\tabcolsep) * \real{0.2060}}
  >{\raggedright\arraybackslash}p{(\columnwidth - 8\tabcolsep) * \real{0.1866}}
  >{\raggedright\arraybackslash}p{(\columnwidth - 8\tabcolsep) * \real{0.2336}}@{}}
\toprule\noalign{}
\begin{minipage}[b]{\linewidth}\raggedright
\textbf{对比对象}
\end{minipage} & \begin{minipage}[b]{\linewidth}\raggedright
\textbf{主要特点}
\end{minipage} & \begin{minipage}[b]{\linewidth}\raggedright
\textbf{优点}
\end{minipage} & \begin{minipage}[b]{\linewidth}\raggedright
\textbf{局限性}
\end{minipage} & \begin{minipage}[b]{\linewidth}\raggedright
\textbf{与本文方法的关系}
\end{minipage} \\
\midrule\noalign{}
\endhead
\bottomrule\noalign{}
\endlastfoot
人工牵引 & 驾驶员和地面人员协同完成牵引作业 &
经验判断灵活，适应真实复杂环境能力强 &
依赖人员经验，标准化和重复精度受人为因素影响 &
本文方法不能直接替代人工，但可为自动化流程提供仿真验证基础 \\
固定路径/开环脚本 & 通过预设速度、转向角和时间完成动作 &
实现简单，适合早期动画和模型验证 &
缺少误差反馈，环境或初始状态变化后适应性差 &
本文在此基础上增加相对位姿误差控制思想 \\
PID控制 & 对速度、角度等单变量进行反馈调节 & 结构简单，工程应用广 &
难以单独处理多变量耦合的精确对接问题 &
可作为底层控制基础，但不足以完整描述对接策略 \\
Pure Pursuit & 基于前视点进行路径跟踪 & 路径跟踪平滑，实现难度适中 &
末端精确姿态对准能力有限 & 适合普通循迹，本文更强调近距离位姿误差修正 \\
Stanley算法 & 基于横向误差和航向误差进行控制 &
适合车辆路径跟踪，横向误差修正能力较好 &
更偏向中高速路径跟踪，低速精细对接针对性不足 &
可作为参考，但本文场景更偏低速对接 \\
MPC控制 & 在约束条件下进行预测优化控制 & 理论效果好，可处理多约束问题 &
建模和调参复杂，实现难度较高 & 后续可作为优化方向，本文阶段暂不采用 \\
SLAM/深度学习方法 & 依赖建图、视觉识别或数据驱动决策 &
自主性和智能化程度高 & 数据、算力和系统复杂度要求高 &
本文保留扩展接口，但现阶段以规则化仿真验证为主 \\
本文方法 & ROS-Gazebo仿真、飞机辅助感知、相对位姿误差控制、4WS执行 &
流程完整、结构清晰、实现难度适中，适合毕业设计验证 &
仿真简化较多，真实环境鲁棒性仍需进一步验证 &
作为自动牵引对接系统的前期仿真与算法验证方案 \\
\end{longtable}

\chapter{ROS-Gazebo仿真系统主要命令与自动牵引作业核心节选}

1.启动命令：

cd~\textasciitilde/tug\_ws

catkin\_make

source~devel/setup.bash

roslaunch~zongzhuang~gazebo.launch

2.控制器加载命令：

rosrun~controller\_manager~spawner~\textbackslash{}

/zongzhuang/front\_left\_wheel\_velocity\_controller~\textbackslash{}

/zongzhuang/front\_right\_wheel\_velocity\_controller~\textbackslash{}

/zongzhuang/back\_left\_wheel\_velocity\_controller~\textbackslash{}

/zongzhuang/back\_right\_wheel\_velocity\_controller

测试脚本运行命令：

chmod~+x~test\_control.py

rosrun~tug\_control~test\_control.py

3.目标点生成

~~~~def~transform\_local\_to\_world(self,~base\_pose,~local\_x,~local\_y,~local\_yaw=0.0):

~~~~~~~~"""

~~~~~~~~将飞机局部坐标系下的点转换到~Gazebo~世界坐标系。

~~~~~~~~"""

~~~~~~~~c~=~math.cos(base\_pose.yaw)

~~~~~~~~s~=~math.sin(base\_pose.yaw)

~~~~~~~~world\_x~=~base\_pose.x~+~c~*~local\_x~-~s~*~local\_y

~~~~~~~~world\_y~=~base\_pose.y~+~s~*~local\_x~+~c~*~local\_y

~~~~~~~~world\_yaw~=~wrap\_to\_pi(base\_pose.yaw~+~local\_yaw)

~~~~~~~~return~Pose2D(world\_x,~world\_y,~world\_yaw)

~~~~def~generate\_task\_points(self,~aircraft\_pose):

~~~~~~~~"""

~~~~~~~~根据飞机位姿生成自动牵引作业所需的关键点：

~~~~~~~~1.~前起落架点；

~~~~~~~~2.~牵引车对接点；

~~~~~~~~3.~预对接点；

~~~~~~~~4.~飞机推出目标点；

~~~~~~~~5.~牵引车撤离点。

~~~~~~~~"""

~~~~~~~~\emph{\#~飞机前起落架点}

~~~~~~~~nose\_gear\_pose~=~self.transform\_local\_to\_world(

~~~~~~~~~~~~aircraft\_pose,

~~~~~~~~~~~~self.nose\_gear\_offset\_x,

~~~~~~~~~~~~self.nose\_gear\_offset\_y,

~~~~~~~~~~~~0.0

~~~~~~~~)

~~~~~~~~\emph{\#~对接点：牵引车前牵引钩与飞机前起落架重合}

~~~~~~~~dock\_pose~=~self.transform\_local\_to\_world(

~~~~~~~~~~~~nose\_gear\_pose,

~~~~~~~~~~~~-self.hitch\_offset\_x,

~~~~~~~~~~~~-self.hitch\_offset\_y,

~~~~~~~~~~~~0.0

~~~~~~~~)

~~~~~~~~\emph{\#~预对接点：沿飞机航向反方向后退一定距离}

~~~~~~~~pre\_dock\_pose~=~self.transform\_local\_to\_world(

~~~~~~~~~~~~dock\_pose,

~~~~~~~~~~~~-self.pre\_dock\_distance,

~~~~~~~~~~~~0.0,

~~~~~~~~~~~~0.0

~~~~~~~~)

~~~~~~~~\emph{\#~推出目标点：飞机沿机身反方向推出}

~~~~~~~~pushback\_pose~=~self.transform\_local\_to\_world(

~~~~~~~~~~~~dock\_pose,

~~~~~~~~~~~~-self.pushback\_distance,

~~~~~~~~~~~~0.0,

~~~~~~~~~~~~0.0

~~~~~~~~)

~~~~~~~~\emph{\#~撤离点：牵引车完成解挂后向侧后方撤离}

~~~~~~~~retreat\_pose~=~self.transform\_local\_to\_world(

~~~~~~~~~~~~pushback\_pose,

~~~~~~~~~~~~self.retreat\_offset\_x,

~~~~~~~~~~~~self.retreat\_offset\_y,

~~~~~~~~~~~~0.0

~~~~~~~~)

~~~~~~~~return~pre\_dock\_pose,~dock\_pose,~pushback\_pose,~retreat\_pose

4.相对位姿误差控制

~~def~compute\_pose\_control(self,~current\_pose,~target\_pose,~max\_speed):

~~~~~~~~"""

~~~~~~~~基于相对位姿误差的车辆控制律。

~~~~~~~~current\_pose：牵引车当前位姿

~~~~~~~~target\_pose：目标位姿

~~~~~~~~max\_speed：当前阶段最大速度

~~~~~~~~输出：

~~~~~~~~speed：车辆线速度

~~~~~~~~curvature：参考路径曲率

~~~~~~~~distance：当前位置到目标点的距离

~~~~~~~~yaw\_error：航向角误差

~~~~~~~~"""

~~~~~~~~dx~=~target\_pose.x~-~current\_pose.x

~~~~~~~~dy~=~target\_pose.y~-~current\_pose.y

~~~~~~~~\emph{\#~将世界坐标误差转换到车辆自身坐标系}

~~~~~~~~c~=~math.cos(current\_pose.yaw)

~~~~~~~~s~=~math.sin(current\_pose.yaw)

~~~~~~~~ex~=~c~*~dx~+~s~*~dy

~~~~~~~~ey~=~-s~*~dx~+~c~*~dy

~~~~~~~~distance~=~math.sqrt(dx~*~dx~+~dy~*~dy)

~~~~~~~~target\_angle\_body~=~math.atan2(ey,~ex)

~~~~~~~~\emph{\#~判断目标点位于车辆前方还是后方}

~~~~~~~~direction~=~1.0

~~~~~~~~alpha~=~target\_angle\_body

~~~~~~~~if~abs(target\_angle\_body)~\textgreater~math.pi~/~2.0:

~~~~~~~~~~~~direction~=~-1.0

~~~~~~~~~~~~alpha~=~wrap\_to\_pi(target\_angle\_body~-~math.pi)

~~~~~~~~yaw\_error~=~wrap\_to\_pi(target\_pose.yaw~-~current\_pose.yaw)

~~~~~~~~\emph{\#~距离误差生成线速度}

~~~~~~~~speed~=~self.k\_rho~*~distance

~~~~~~~~speed~=~clamp(speed,~self.min\_motion\_speed,~max\_speed)

~~~~~~~~speed~*=~direction

~~~~~~~~\emph{\#~横向误差与航向误差共同生成角速度}

~~~~~~~~if~distance~\textgreater~0.8:

~~~~~~~~~~~~omega~=~self.k\_alpha~*~alpha

~~~~~~~~else:

~~~~~~~~~~~~omega~=~self.k\_alpha~*~alpha~+~self.k\_yaw~*~yaw\_error

~~~~~~~~\emph{\#~倒车时角速度方向修正}

~~~~~~~~if~direction~\textless~0.0:

~~~~~~~~~~~~omega~=~-omega

~~~~~~~~\emph{\#~角速度转换为参考路径曲率}

~~~~~~~~if~abs(speed)~\textless~1e-5:

~~~~~~~~~~~~curvature~=~0.0

~~~~~~~~else:

~~~~~~~~~~~~curvature~=~omega~/~speed

~~~~~~~~return~speed,~curvature,~distance,~abs(yaw\_error)

5.自动牵引作业状态机主流程

~~~~def~run\_task\_state\_machine(self,~aircraft\_pose,~current\_tug\_pose):

~~~~~~~~"""

~~~~~~~~自动牵引作业状态机主逻辑。

~~~~~~~~该函数体现从目标识别、自动对接、推出到撤离的完整流程。

~~~~~~~~"""

~~~~~~~~pre\_dock\_pose,~dock\_pose,~pushback\_pose,~retreat\_pose~=~\textbackslash{}

~~~~~~~~~~~~self.generate\_task\_points(aircraft\_pose)

~~~~~~~~\emph{\#~搜索并确认飞机目标}

~~~~~~~~if~self.state~==~TaskState.SEARCH\_TARGET:

~~~~~~~~~~~~target\_valid~=~aircraft\_pose~is~not~None

~~~~~~~~~~~~if~target\_valid:

~~~~~~~~~~~~~~~~self.state~=~TaskState.APPROACH\_PRE\_DOCK

~~~~~~~~~~~~return~0.0,~0.0,~0.0

~~~~~~~~\emph{\#~前往预对接点}

~~~~~~~~elif~self.state~==~TaskState.APPROACH\_PRE\_DOCK:

~~~~~~~~~~~~speed,~front\_delta,~back\_delta,~reached~=~self.go\_to\_pose\_step(

~~~~~~~~~~~~~~~~current\_tug\_pose,

~~~~~~~~~~~~~~~~pre\_dock\_pose,

~~~~~~~~~~~~~~~~self.max\_approach\_speed,

~~~~~~~~~~~~~~~~"front\_only"

~~~~~~~~~~~~)

~~~~~~~~~~~~if~reached:

~~~~~~~~~~~~~~~~self.state~=~TaskState.PRECISE\_DOCK

~~~~~~~~~~~~return~speed,~front\_delta,~back\_delta

~~~~~~~~\emph{\#~低速精确对接}

~~~~~~~~elif~self.state~==~TaskState.PRECISE\_DOCK:

~~~~~~~~~~~~speed,~front\_delta,~back\_delta,~reached~=~self.go\_to\_pose\_step(

~~~~~~~~~~~~~~~~current\_tug\_pose,

~~~~~~~~~~~~~~~~dock\_pose,

~~~~~~~~~~~~~~~~self.max\_dock\_speed,

~~~~~~~~~~~~~~~~"counter\_phase"

~~~~~~~~~~~~)

~~~~~~~~~~~~if~reached:

~~~~~~~~~~~~~~~~self.state~=~TaskState.ATTACH

~~~~~~~~~~~~return~speed,~front\_delta,~back\_delta

~~~~~~~~\emph{\#~建立牵引连接}

~~~~~~~~elif~self.state~==~TaskState.ATTACH:

~~~~~~~~~~~~\emph{\#~工程实现中，此处可建立虚拟牵引约束或调用连接装置模型}

~~~~~~~~~~~~self.state~=~TaskState.PUSHBACK

~~~~~~~~~~~~return~0.0,~0.0,~0.0

~~~~~~~~\emph{\#~飞机推出}

~~~~~~~~elif~self.state~==~TaskState.PUSHBACK:

~~~~~~~~~~~~speed,~front\_delta,~back\_delta,~reached~=~self.go\_to\_pose\_step(

~~~~~~~~~~~~~~~~current\_tug\_pose,

~~~~~~~~~~~~~~~~pushback\_pose,

~~~~~~~~~~~~~~~~self.max\_pushback\_speed,

~~~~~~~~~~~~~~~~"counter\_phase"

~~~~~~~~~~~~)

~~~~~~~~~~~~if~reached:

~~~~~~~~~~~~~~~~self.state~=~TaskState.RELEASE

~~~~~~~~~~~~return~speed,~front\_delta,~back\_delta

~~~~~~~~\emph{\#~解除牵引连接}

~~~~~~~~elif~self.state~==~TaskState.RELEASE:

~~~~~~~~~~~~\emph{\#~工程实现中，此处解除虚拟牵引约束或机械连接}

~~~~~~~~~~~~self.state~=~TaskState.RETREAT

~~~~~~~~~~~~return~0.0,~0.0,~0.0

~~~~~~~~\emph{\#~车辆撤离}

~~~~~~~~elif~self.state~==~TaskState.RETREAT:

~~~~~~~~~~~~speed,~front\_delta,~back\_delta,~reached~=~self.go\_to\_pose\_step(

~~~~~~~~~~~~~~~~current\_tug\_pose,

~~~~~~~~~~~~~~~~retreat\_pose,

~~~~~~~~~~~~~~~~self.max\_retreat\_speed,

~~~~~~~~~~~~~~~~"front\_only"

~~~~~~~~~~~~)

~~~~~~~~~~~~if~reached:

~~~~~~~~~~~~~~~~self.state~=~TaskState.DONE

~~~~~~~~~~~~return~speed,~front\_delta,~back\_delta

~~~~~~~~\emph{\#~任务完成}

~~~~~~~~elif~self.state~==~TaskState.DONE:

~~~~~~~~~~~~return~0.0,~0.0,~0.0

~~~~~~~~\emph{\#~异常中止}

~~~~~~~~else:

~~~~~~~~~~~~self.state~=~TaskState.ABORT

~~~~~~~~~~~~return~0.0,~0.0,~0.0
